% Proposed isolated B165, new checkpoint34 source. No earlier bytes changed.
\subsection*{The entire first cohomology of the actual deletion kernel}
\statement{B.165}{Proved}{The whole deletion-kernel cohomology is the completed source coset space, and the filling-kernel cohomology maps onto it with a continuous pointed section}
\label{stmt:B165}
Retain the actual generic ordinary group of B.123/B.135 and the
\emph{filled} composite of B.139:
\begin{equation}\label{eq:B165-filled}
 G_{\mathrm{gen}}=F_{11}\rtimes S
 \longrightarrow \Gamma_0=G_{\mathrm{gen}}/N
 \longrightarrow\Gamma_a,\qquad q_a=\text{the composite},
 \quad S=F(m,v).
\end{equation}
Here \(N\) is the whole actual node normal closure, with all forty
ordinary seeds and every signed source translate. Put
\[
 \begin{gathered}
 K_a=\ker q_a,\qquad
 J=\Gamma_a^\delta
   =\langle A,m,v\mid[A,m]=1\rangle\cong\mathbb Z^2*\mathbb Z,
 \qquad P=q_a^{-1}(J),\\
 \mathcal K_a=\overline{K_a},\qquad
 \mathcal N=\overline N,\qquad
 \mathcal P=\overline P,\qquad
 \mathcal R=\overline R .
 \end{gathered}
\]
Every closure is taken in the full completion of the actual generic
group. The section \(\lambda:J\to G_{\mathrm{gen}}\) sends
\(A\) to \(aa^*\), where \(a^*=\tau a\tau\), and fixes \(m,v\);
its completed embedding is the actual one of B.159/B.162.
Let \(\Delta=\langle\delta\rangle\cong C_2\), with
\(\delta=\operatorname{Inn}(\tau)\). Then
\begin{equation}\label{eq:B165-entire}
 \widehat J/\widehat S
 \xrightarrow{\ \sim\ } H^1(\Delta,\mathcal K_a),
 \qquad
 j\widehat S\longmapsto[X_j],
 \quad X_j=\delta\widehat\lambda(j)\widehat\lambda(j)^{-1},
\end{equation}
is a pointed homeomorphism. In particular the map
\[
 H^1(\Delta,\mathcal K_a)\longrightarrow H^1(\Delta,\mathcal P)
\]
is constant at the basepoint. Moreover
\begin{equation}\label{eq:B165-section}
 H^1(\Delta,\mathcal N)\longrightarrow H^1(\Delta,\mathcal K_a)
\end{equation}
is surjective and has the continuous pointed section supplied by
\(H^1(\Delta,\mathcal R)\). These are statements about compact
pointed spaces, not asserted group splittings. The kernel fiber of
\eqref{eq:B165-section} is not identified with the contact image.

\begin{proof}
\emph{The actual completed preimage.}
The filled map preserves the split source projection to \(S\).
Right exactness of full profinite completion gives
\(\ker\widehat q_a=\mathcal K_a\).
B.139 and B.154 give the full induced topology of \(J\) and the
actual equality
\[
 (\widehat{\Gamma_a})^\delta=\widehat J.
\]
The completed section therefore proves
\begin{equation}\label{eq:B165-preimage}
 \widehat q_a^{-1}(\widehat J)
   =\mathcal K_a\rtimes\widehat\lambda(\widehat J)
   =\mathcal P.
\end{equation}
Indeed the displayed product is closed and contains \(P\);
conversely both of its factors lie in \(\overline P\).
This proves the precise preimage equality without identifying
\(\mathcal K_a\) with the abstract completion of an infinite-index
subgroup.

\emph{Every deletion-kernel cocycle is a preimage coboundary.}
Let \(c\in\mathcal K_a\) satisfy \(c\delta(c)=1\).
Because the filled map retains the source projection, \(\mathcal K_a\)
has trivial source coordinate and lies in the actual completed
ordinary free fiber \(\widehat F_{11}\).
The full-topology and orientation comparisons of B.123/B.135 identify
its orbifold fiber with the finite-factor free profinite product
\begin{equation}\label{eq:B165-fiber}
 \widehat E_0
 =\widehat F(a,b,d,e,f)\amalg C_2(C)\amalg C_2(\tau).
\end{equation}
The original orientation has ordinary kernel \(\widehat F_{11}\).
Thus \(c\tau\) is a nonidentity involution in \(\widehat E_0\):
its square is \(c\delta(c)=1\), and its orientation is odd.
The prime-order torsion theorem of
\citet[Theorem A\(^{\prime}\)]{A_HR1985} makes it conjugate
to \(C\) or \(\tau\). The free profinite factor is torsion-free.
Only this finite-factor theorem is used; no infinite-index Kurosh
comparison is assumed.

The \(C\) alternative is excluded in the \emph{actual filled target}.
Extend \(q_a\) to the marked orbifold map, fixing the selected
\(\tau\). Its target has the continuous character
\[
 C\longmapsto(1,0),\qquad \tau\longmapsto(0,1),
 \qquad a,f,m,v\longmapsto(0,0)
 \quad\text{in }C_2^2.
\]
It factors through the coarse \(Q_f\rtimes S\) quotient of B.139:
the entire signed source actions preserve the two characters and
the surviving node commutators vanish. In particular it is a
character of the filled group, not merely of the pre-node deletion.
Since \(q_a(c\tau)=\tau\), the two unequal character values rule
out conjugacy to \(C\). Consequently
\[
                         c\tau=g\tau g^{-1}
                         \qquad(g\in\widehat E_0).
\]
If \(g\) is odd, replace it by \(g\tau\); this preserves the conjugate
and makes it even. Thus \(g\in\widehat F_{11}\).
Its filled image is ordinary and centralizes \(\tau\), so belongs
to \(\widehat J\) by the completed one-family fixed theorem.
Equation \eqref{eq:B165-preimage} now gives \(g\in\mathcal P\).
Finally
\[
 c=g\delta(g)^{-1}=(g^{-1})^{-1}\delta(g^{-1}),
\]
which is a \(\mathcal P\)-coboundary in the stated convention.
This proves that the map to \(H^1(\Delta,\mathcal P)\) is constant.

\emph{Entire target and its topology.}
B.163 already proves that the map in \eqref{eq:B165-entire}
is injective with image exactly
\[
 \ker\!\left[
 H^1(\Delta,\mathcal K_a)\longrightarrow H^1(\Delta,\mathcal P)
 \right].
\]
The preceding argument makes this kernel the entire target.
The map is continuous. The cocycle set in a compact group is
closed, and the compact group action
\(c\mapsto b^{-1}c\delta(b)\) has closed orbit relation;
therefore its pointed cohomology space is compact Hausdorff.
The resulting continuous bijection from
\(\widehat J/\widehat S\) is a homeomorphism.

Finally B.161--B.163 give
\(\mathcal R\le\mathcal N\le\mathcal K_a\) and identify
\(H^1(\Delta,\mathcal R)\to H^1(\Delta,\mathcal K_a)\)
with this same homeomorphism. Its inverse followed by the inclusion
map to \(H^1(\Delta,\mathcal N)\) is the claimed continuous pointed
section of \eqref{eq:B165-section}.
It does not determine the basepoint fiber. In particular B.164's
known \(D_s\) contact is a nonbase class in that fiber; whole-map
injectivity remains false.
\end{proof}

\begingroup\footnotesize\raggedright
The complete prior proof and exact source bindings are retained in
\path{work_CE/entire_deleted_kernel_cohomology33.txt} and
\path{work_CE/pre_node_input_binding33.txt}.
Its complete independent review is
\path{work_DE/review_entire_deleted_kernel_cohomology33.md}.
This strengthens the target assertion of B.163. It neither proves
the remaining contact-image equality of B.164 nor the separate
global Main B comparison.
\par\endgroup
